\section{Ontology}
\begin{frame}{Ontologie}
		\begin{figure}
			\includegraphics[width=.9\linewidth]{T-Box}
			\caption{Ontology T-Box}
		\end{figure}
\end{frame}


\begin{frame}[fragile]{Quantum Ontology}
	\begin{minipage}{.35\linewidth}
		\begin{block}{}
			\centering
			Traduzione in Prolog
		\end{block}
		\begin{figure}[h]
			\begin{minipage}{\linewidth}
				\hrule
				\begin{minted}{Prolog}

class(animal).
class(mammal).

subclass(mammal, animal).
subclass(seal, mammal).
				\end{minted}
				\hrule
			\end{minipage}
			\caption{Classes and subclasses}
		\end{figure}
\vspace{-.7cm}
		\begin{figure}[h]
			\begin{minipage}{\linewidth}
				\hrule
				\begin{minted}{Prolog}

subClass(C1, C2) :-
	subclass(C1, C2).
	
subClass(C1, C2) :-
	subclass(C1, C3),
	subClass(C3, C2).
				\end{minted}
				\hrule
			\end{minipage}
			\caption{Subclasses recursive definition}
		\end{figure}
	\end{minipage}
	\hfill
	\pause
	\begin{minipage}{.55\linewidth}
		\begin{block}{}
			\centering
			Unfolding
		\end{block}
		\begin{figure}[h]
			\hrule
			\begin{minipage}{.45\linewidth}
				\begin{minted}{Prolog}
					
subClass0(C0, C1) :-
	subclass(C0, C1).
	
subClass1(C0, C1) :-
	subClass0(C0, C2),
	subclass(C2, C1).
	
				\end{minted}
			\end{minipage}
			\hfill
			\begin{minipage}{.45\linewidth}
				\begin{minted}{Prolog}
subClass2(C0, C1) :-
	subClass1(C0, C2),
	subclass(C2, C1)
	
subClass3(C0, C1) :-
	subClass2(C0, C2),
	subclass(C2, C1)
				\end{minted}
			\end{minipage}
			\hrule
			\caption{Subclasses rule unfolded}
		\end{figure}
		\vspace{-.5cm}
		\begin{figure}[h]
			\begin{minipage}{\linewidth}
				\hrule
				\begin{minted}{Python}
					
rule_0 = 'subClass0(C0,C1) :- \n\tsubclass(C0, C1).\n'
for i in range(1, n_unfold):
	rule = 'subClass'+ str(i) + '(C0,C2) :-\n'        +\
		     '\tsubClass' + str(i-1) + '(C0,C3),\n'     +\
			   '\tsubclass(C3, C2).\n'
				\end{minted}
				\hrule
			\end{minipage}
			\caption{Unfolding script}
		\end{figure}
	\end{minipage}
\end{frame}

\begin{frame}[fragile]{Dimensione del problema}
	\hfill
		\begin{minipage}{.43\linewidth}
			\begin{block}{}
				\centering
				Query
			\end{block}
			\begin{figure}[h]
				\begin{minipage}{\linewidth}
					\hrule
					\begin{enumerate}
						\vspace{.1cm}
						\item \mintinline{Prolog}|?- class(animal).|\vspace{.1cm}
						\item \mintinline{Prolog}|?- class(X).|\vspace{.1cm}
						\item \mintinline{Prolog}|?- subClass(seal, X).|\vspace{.1cm}
						\item \mintinline{Prolog}|?- isA(sardy, X).|\vspace{.1cm}
						\item \mintinline{Prolog}|?- error(X, Y, Z).|\vspace{.1cm}
						\item \mintinline{Prolog}|?- hasProperty(I1, ancestor, reiny_c).|
					\end{enumerate}
					\hrule
				\end{minipage}
				\caption{Example Query}
			\end{figure}
		\end{minipage}
		\hfill
		\begin{minipage}{.48\linewidth}
			\begin{figure}
				\includegraphics[width=.9\linewidth]{dimension}
				\caption{Problem Dimension}
			\end{figure}
		\end{minipage}
	\end{frame}

\begin{frame}[fragile]{Dimensione del problema}{Impatto dell'unfolding}
	\hfill
	\begin{minipage}{.4\linewidth}
		\begin{block}{}
			\centering
			Knowledge Base
		\end{block}
		\begin{figure}[h]
			\begin{minipage}{\linewidth}
				\hrule
				\begin{minted}{Prolog}
					
class(animal).
class(mammal).

subclass(mammal, animal).
				\end{minted}
				\hrule
			\end{minipage}
			\caption{Facts}
		\end{figure}
		\begin{figure}[h]
			\begin{minipage}{\linewidth}
				\hrule
				\begin{minted}{Prolog}
					
?- subClass(mammal, animal)
				\end{minted}
				\hrule
			\end{minipage}
			\caption{Query}
		\end{figure}
	\end{minipage}
	\hfill
	\begin{minipage}{.48\linewidth}
		\begin{figure}
			\includegraphics<1>[width=.9\linewidth]{unfold}
			\includegraphics<2>[width=.9\linewidth]{unfold_fit}
			\caption{Problem Dimension}
		\end{figure}
	\end{minipage}
\end{frame}